\paragraph{The fermion form of the doubleton member at fixed \(\gamma>0\) (conditional on the class of operators; bounded result).}
Keep the setting of the preceding paragraph and take the member with \(T^a=M^{a4}\), \(a=0,\dots,3\), on a doubleton representation \(\mathcal H_n\) of \(\mathfrak{so}(4,2)\)~\cite{MantaSteinacker2025}, with \(W\) timelike, \(c>0\) and \(\gamma=3c\). The fermion fluctuations are taken to be Hilbert--Schmidt operators on \(\mathcal H_n\), with the inner product \(\operatorname{Tr}\psi^\dagger\psi\), times the spinor index, and the form is taken on the vectors that are smooth for the adjoint action of \(\mathrm{Spin}(4,1)\). Both are premises, since whether a kernel is normalisable depends on them. Under \(\mathfrak{so}(3,1)\times\mathfrak{so}(6)\) the chiral spinor of \(\mathrm{SO}(9,1)\) consists of four copies of \((\tfrac12,0)\) and four of \((0,\tfrac12)\), so the form is non-degenerate exactly when the Weyl operators \(\mathcal D_\pm\psi=\sigma_\pm^a[\bar T_a,\psi]\) have trivial kernels, where \(\sigma_\pm^a\) are the blocks of \(\Gamma^a\) between the two chiralities of \(\mathfrak{so}(3,1)\), so that \(\sigma_\mp^a\sigma_\pm^b+\sigma_\mp^b\sigma_\pm^a=2\eta^{ab}\). A member with another frame \(L\) is carried to this one by an isometry of \(h_0\), whose lift acts invertibly on the spinor and maps each chirality to one of the two.

\emph{Algebra.} On one chirality the identity of the preceding paragraph reads \(c^{-1}\mathcal D_-\mathcal D_+=C_5-C_4-2X\), where \(C_5\) and \(C_4\) are now the Casimirs of \(\mathfrak{so}(4,1)\) and \(\mathfrak{so}(3,1)\) formed with \(\eta\), \(X=\tfrac12(C_4^{\mathrm{tot}}-C_4-\tfrac32)=\vec\sigma\cdot\vec J_{\mathrm L}\), and \(\vec J_{\mathrm L}=\tfrac12(\vec J\pm\mathrm i\vec K)\) is formed from the rotations \(\vec J\) and boosts \(\vec K\) of \(\mathfrak{so}(3,1)\) in the adjoint action; the other chirality has \(\vec J_{\mathrm R}\) in its place. Then \(X^2+X=\vec J_{\mathrm L}^2\) and \(4\vec J_{\mathrm L}^2=C_4\pm2\mathrm iP\), with \(P=\vec J\cdot\vec K\) and the sign fixed by the chirality (checked on explicit finite-dimensional representations of the real form \(\mathfrak{so}(4,1)\)). Let \(\psi\) be in the kernel and put \(A=C_5-C_4\). Then \(2X\psi=A\psi\), so \(4X^2\psi=A^2\psi\), and therefore \(B\psi=\mp2\mathrm iP\psi\) with \(B=1+C_4-(A+1)^2\). On smooth vectors \(B\) and \(P\) are symmetric and commute, so \(\|B\psi\|^2=\mp2\mathrm i\langle B\psi,P\psi\rangle\) with \(\langle B\psi,P\psi\rangle\) real, and hence \(B\psi=P\psi=0\). For the compact members \(\vec J_{\mathrm L}\) is Hermitian and no such condition arises; the factor \(\mathrm i\) comes from the real form.

\emph{Spectral condition.} The Casimirs act as real scalars on the irreducible components, \(C_5\) for \(\mathrm{Spin}(4,1)\) and \(C_4\), \(P\) for \(\mathrm{Spin}(3,1)\), and so extend to strongly commuting self-adjoint operators. Hence \(\psi\) lies in the spectral subspace of \((C_4,P)\) for the set \(P=0\), \(1+C_4\ge0\). On the unitary principal series of \(\mathrm{Spin}(3,1)\cong\mathrm{SL}(2,\mathbb C)\), labelled by \(k\in\tfrac12\mathbb Z\) and \(p\in\mathbb R\), one has \(C_4=\vec J^2-\vec K^2=k^2-p^2-1\) and \(P=\pm kp\) (Eq.~(297) of~\cite{MartinDussaud2019}; rederived in the realisation by functions on \(\mathbb C\)). Within the principal series this set therefore consists of the countably many representations with \(p=0\).

\emph{Restriction to \(\mathrm{Spin}(3,1)\).} The \(M^{\mu\nu}\) and \(M^{\mu4}\pm M^{\mu5}\) span Poincar\'e subalgebras, and the ladder representations restrict to the Poincar\'e group as irreducible massless representations of fixed helicity~\cite{MackTodorov1969}, as stated in~\cite{Todorov2019}; for \(n>0\) this uses the identification of the doubleton \(\mathcal H_n\) with a ladder representation, which is taken from the literature on doubletons and not independently verified here. As a representation of \(\mathrm{SL}(2,\mathbb C)\), \(\mathcal H_n\) is then a space of sections over the light cone, and the Hilbert--Schmidt operators, \(\mathcal H_n\otimes\bar{\mathcal H}_n\), are sections over pairs of null vectors. The pairs that are not collinear have full measure; they form orbits labelled by their scalar product, with stabiliser \(\mathrm U(1)\), which is compact. Each orbit therefore carries a representation induced from \(\mathrm U(1)\), which is contained in the regular representation, so the adjoint action restricted to \(\mathrm{Spin}(3,1)\) is contained in a multiple of the regular representation. The Plancherel measure is supported on the principal series and is the counting measure in the discrete label times a polynomial density times Lebesgue measure in the continuous one (Theorem~7.1 of~\cite{FioresiYuncken2024}). The set above therefore has measure zero, and \(\psi=0\). For \(n=0\) this also follows without the Poincar\'e group. On \(\mathcal H_0\) one has \(C_4=-1-(M^{45})^2\) and \(P=0\) (Eqs.~(11) and~(12) of~\cite{MantaSteinacker2025}), so, by the classification of the unitary dual in~\cite{MartinDussaud2019}, \(\mathcal H_0\) restricts to principal series with \(k=0\), and a tensor product of principal series is a direct integral of principal series over Lebesgue measure in \(p\) (Eq.~(308) of~\cite{MartinDussaud2019}).

So the fermion form of the doubleton member has no normalisable kernel in this class, for \(n=0\) and, given that identification, for every \(n\). The hypothesis of the conditional statement of the preceding paragraph thus holds for these members in that sense. The decomposition of the operators under \(\mathfrak{so}(4,1)\) is not needed. The class matters: among bounded operators the identity, times any constant spinor, lies in the kernel, as the trace mode does at finite \(N\) before it is removed, and the kernel in that class is not computed. A trivial kernel does not make the operator invertible with a bounded inverse. Whether zero lies in its continuous spectrum, how the Pfaffian at \(N=\infty\) is then defined, and whether it vanishes are not decided, so the separation at one loop stated conditionally in the preceding paragraph is not established. Kernel vectors that are not smooth, members on other unitary representations of \(\mathfrak{so}(4,1)\), and products with a factor \(K\) are not treated. These statements concern the operator at the member, for \(N\to\infty\) taken first at fixed \(\gamma>0\), and involve no \(\epsilon\), which Definition~(A) puts on the bosons only. The family, \(h_0\), the mass and the fermion content remain supplied, so this removes one obstruction to the conditional route of the preceding paragraph without deriving the signature, and Problem~E stays open with Assumption~\ref{ass:lorentzian} retained.
